\section{Beurteilung des Verlaufs des Aktienkurses}
\label{sec:kurs}

Grundlage sind die Nasdaq-Schlusskurse aus Abbildung~\ref{fig:kursverlauf}
und Tabelle~\ref{tab:kursextrema} sowie die Ergebnisse der
Teile~\ref{sec:organisch} und \ref{sec:accession}.

\begin{figure}[htbp]
\centering
\begin{tikzpicture}
\begin{axis}[kursachse]
\addplot+[gray!55, thick, mark=none] table[col sep=comma, x=datum,
  y=kurs] {data/kursverlauf.csv};
\addplot+[kurspunkte] table[col sep=comma, x=datum,
  y=kurs, meta=art] {data/kursverlauf.csv};
\end{axis}
\end{tikzpicture}
\caption{Kursverlauf der Brown-\&-Brown-Aktie vom Jahresschluss 2020 bis zum
27.~August 2026 \"uber achtzehn St\"utzpunkte. Jeder Punkt steht an seinem
tats\"achlichen Handelstag: {\color{kurshoch}$\blacktriangle$} Jahreshoch,
{\color{kurstief}$\blacktriangledown$} Jahrestief,
{\color{kursschluss}$\bullet$} Jahresschluss.}
\label{fig:kursverlauf}
\end{figure}

\begin{table}[htbp]
\centering
\caption{Jahresschluss-, H\"ochst- und Tiefstkurse in USD. Das Jahr 2026
reicht bis zum 27.~August.}
\label{tab:kursextrema}
\begin{tabular}{lrrrl}
\toprule
Jahr & Schluss & H\"ochstkurs & Tiefstkurs & Ver\"anderung\\
\midrule
2021 & \Kurswert{\KursSchlussZFI} & \Kurswert{\KursHochZFI} & \Kurswert{\KursTiefZFI} & --\\
2022 & \Kurswert{\KursSchlussZFII} & \Kurswert{\KursHochZFII} & \Kurswert{\KursTiefZFII} & \Pct{\fpeval{round(\KursDeltaZweiZwei,1)}}\\
2023 & \Kurswert{\KursSchlussZFIII} & \Kurswert{\KursHochZFIII} & \Kurswert{\KursTiefZFIII} & \Pct{\fpeval{round(\KursDeltaZweiDrei,1)}}\\
2024 & \Kurswert{\KursSchlussZFIV} & \Kurswert{\KursHochZFIV} & \Kurswert{\KursTiefZFIV} & \Pct{\fpeval{round(\KursDeltaZweiVier,1)}}\\
2025 & \Kurswert{\KursSchlussZFV} & \Kurswert{\KursHochZFV} & \Kurswert{\KursTiefZFV} & \Pct{\fpeval{round(\KursDeltaZweiFuenf,1)}}\\
2026 & \Kurswert{\KursSchlussZFVI} & \Kurswert{\KursHochZFVI} & \Kurswert{\KursTiefZFVI} & \Pct{\fpeval{round(\KursDeltaLaufend,1)}}\\
\bottomrule
\end{tabular}
\end{table}

\FloatBarrier

% FAKTENBASIS 4
%  - Jahresschlusskurse und Extrema: Tabelle \ref{tab:kursextrema}
%  - Organisches Wachstum: 2024 \OrganischWachstumVJ %, 2025
%    \OrganischWachstum %, H1 2026 \HJOrganisch % (Abschnitt \ref{sec:organisch})
%  - Umsatz 2025 +\UmsatzDelta %, bereinigtes EBITDAC +\EBITDACberDelta %,
%    beide Rekordwerte (Tabelle \ref{tab:kennzahlen})
%  - Ergebnis je Aktie \EPSdilDelta % - faellt trotz Ergebniswachstum
%  - Hoechstkurs 2025 \KursHochZFV USD am \Datum{\KurstagHochZFV}, also VOR der
%    Ankuendigung der Finanzierung im Juni 2025
%  - Tiefstkurs 2026 \KursTiefZFVI USD am \Datum{\KurstagTiefZFVI}
%  - Nettoverschuldung \NettoverschuldungVJ -> \Nettoverschuldung Mio. USD,
%    Verschuldungsgrad \VerschuldungsgradVJ -> \Verschuldungsgrad
%  - Kaufpreis der Zukaeufe \KaufpreisJeUmsatz-facher Jahresumsatz
%  - Keine Analystenratings oder Kursziele belegbar (Abschnitt \ref{sec:luecken})

\bk{Der Kursverlauf l\"asst sich aus den vorstehenden Analysen erkl\"aren,
allerdings nicht anhand der Gewinn- und Verlustrechnung, sondern anhand des
organischen Wachstums. Es ist f\"ur einen Makler der Fr\"uhindikator der
k\"unftigen Ertragskraft, w\"ahrend der berichtete Umsatz auch das
wiedergibt, was zugekauft wurde. Dem Kursanstieg um
\Pct{\KursDeltaZweiVier} im Jahr 2024 steht ein organisches Wachstum von
\Pct{\OrganischWachstumVJ} gegen\"uber, die h\"ochste der hier betrachteten
Raten. Umgekehrt verlor die Aktie 2025 \PctBetrag{\KursDeltaZweiFuenf},
obwohl Umsatz (\Pct{\UmsatzDelta}) und bereinigtes EBITDAC
(\Pct{\EBITDACberDelta}) Rekordwerte erreichten. Erkl\"arungskr\"aftig ist
allein das organische Wachstum, das im selben Jahr von
\Pct{\OrganischWachstumVJ} auf \Pct{\OrganischWachstum} fiel.

Wie fr\"uh der Markt reagierte, zeigt der unterj\"ahrige Verlauf. Das
Jahreshoch 2025 von \USDje{\KursHochZFV} lag am \Datum{\KurstagHochZFV} und damit
vor der Ank\"undigung der \"Ubernahme und ihrer Finanzierung im Juni; das
Jahrestief von \USDje{\KursTiefZFV} folgte am \Datum{\KurstagTiefZFV}, also nach
der ersten vollen Berichtsperiode mit Accession. Zwischen beiden Punkten
liegt der Wechsel der Frage, die der Markt stellt: nicht mehr, wie schnell
das Unternehmen w\"achst, sondern zu welchem Preis.

Drei Gr\"o{\ss}en beantworten diese Frage ung\"unstig. Die Nettoverschuldung
stieg von \MioUSD{\NettoverschuldungVJ} auf \MioUSD{\Nettoverschuldung} und
damit von \Faktor{\VerschuldungsgradVJ} auf \Faktor{\Verschuldungsgrad} des
bereinigten EBITDAC. Die Zahl der Aktien wuchs um \Pct{\AktienZuwachs}, so
dass das Ergebnis je Aktie trotz eines um \Pct{\ErgebnisDelta} h\"oheren
Konzernergebnisses um \PctBetrag{\EPSdilDelta} zur\"uckging. Und der
Kaufpreis der Zuk\"aufe entspricht dem \Faktor{\KaufpreisJeUmsatz}-fachen
ihres Jahresumsatzes, w\"ahrend das eigene Unternehmen zum Jahresschluss
2025 mit dem \Faktor{\fpeval{round(\Marktkapitalisierung/\UmsatzGesamt,2)}}-fachen
seines Umsatzes bewertet wurde. Wer teurer kauft, als er selbst bewertet
ist, muss die Differenz durch Wachstum oder Kostensenkung verdienen; genau
das Wachstum ist im selben Jahr eingebrochen.

Der Verlauf 2026 best\"atigt diese Lesart. Das organische Wachstum wurde im
ersten Halbjahr mit \Pct{\HJOrganisch} negativ, und der Kurs fiel bis zum
\Datum{\KurstagTiefZFVI} auf \USDje{\KursTiefZFVI}, den tiefsten Stand seit 2022.
Der aktuelle Kurs von \USDje{\KursSchlussZFVI} liegt
\PctBetrag{\KursAbstandHoch} unter dem H\"ochstkurs von April 2025. Nicht
aus den Ergebnissen erkl\"arbar bleibt, wie weit der Markt die
Integrationsrisiken bereits eingepreist hat: Ratings und Kursziele der
achtzehn abdeckenden H\"auser sind nicht ver\"offentlicht
(Abschnitt~\ref{sec:luecken}), so dass die Erwartung des Marktes nur \"uber
den Kurs selbst ablesbar ist.}
